\documentclass[10pt,a4paper]{amsart}
\usepackage[margin=19mm]{geometry}
\usepackage{amsmath,amssymb,mathtools}
\usepackage[scr=rsfs,cal=euler]{mathalfa}
\usepackage{xfrac}
\usepackage[unicode,hidelinks,bookmarks=false]{hyperref}
\newcommand{\ie}{i.e.\ }
\newcommand{\cf}{cf.\ }
\newcommand{\resp}{respectively\ }
\newcommand{\Subsection}{\subsection}
\providecommand{\nobreakdash}{\nobreak\hbox{-}}
\setlength{\emergencystretch}{2em}
\multlinegap0pt
\usepackage{amssymb}
\usepackage{xcolor}
\usepackage{enumerate}
\usepackage{float}
\usepackage[shortlabels]{enumitem}
\usepackage{tikz}
\usepackage[normalem]{ulem}
\newtheorem{lemma}{Lemma}
\title{On difference-differential Lax pairs and integrals of Painlev\'e equations in finite characteristic}
\author{Nalini Joshi, Tomas Lasic Latimer, Pieter Roffelsen}
\date{Source: arXiv:2607.06980v1 (CC BY 4.0)}
\begin{document}
\maketitle

\noindent\emph{Source and licence.} On difference-differential Lax pairs and integrals of Painlev\'e equations in finite characteristic. Version arXiv:2607.06980v1 (CC BY 4.0), distributed by arXiv under the Creative Commons Attribution 4.0 International licence, http://creativecommons.org/licenses/by/4.0/, as recorded in the arXiv licence metadata for this exact version and shown on the abstract page https://arxiv.org/abs/2607.06980v1. Full licence notice: \texttt{licenses/arxiv-2607.06980-CC-BY-4.0.txt}.\par\smallskip

\noindent\textit{Editorial scope.} Counted unit: the article's lemma lem:trace (source0.tex lines 371--380). The article's complete original proof of it is delivered with it (source0.tex lines 1711--1779, one \texttt{proof} environment of 69 lines, which the article places away from the statement). No mathematics is written, repaired or re-derived: the statement and the proof are the article's own lines, in the article's own notation, and no retained line is substituted or re-flowed.  Retained uncounted as the source-local context the proof needs: lines 335--337; lines 362--364; lines 1630--1641; lines 1679--1681; lines 1701--1703.  Nothing was omitted from any retained span, so the package carries no omission record.  No retained line contains a citation, so the wrapper carries no bibliography and no bibliography entry.  No retained line contains a figure environment, an image-inclusion command or drawing code, so no image is delivered and the package declares no asset dependency.  Editorial formatting only, in addition: an AMS \texttt{amsart} preamble that restates the article's own declarations and macros used by the retained lines, namely the environments \texttt{lemma} on separate counters, together with the standard packages those lines need.  The retained environments print in this document as Lemma 1 in order of appearance, whereas the article prints the same items with the numbering of its own full document.  The complete native source of the article as distributed in the arXiv e-print for this exact version is bundled unchanged as \texttt{sources/204674/source0.tex}.
\begin{equation}\label{eq:Mp def}
    M_p(z)=A(z+p-1)\cdot A(z+p-2)\cdot\ldots\cdot A(z+1)\cdot A(z),
\end{equation}

\begin{equation}\label{eq:notationps_reform}
 Z=z^p-z,\quad  T=t^p,\quad \mathcal{A}_k=\alpha_k^p-\alpha_k\quad (k\in\mathbb{Z}).
\end{equation}

\begin{lemma}\label{lemma:z+n}
For any odd prime $p$ and integer $0\leq N\leq 2p-1$, we have the following identities in  $\mathbb{F}_p[z]$,
\begin{equation}\label{eq:sumidentity}
\sum_{n=0}^{p-1}(z+n)^N\, = \begin{cases}
0 &\text{if $0\leq N < p-1$},\\
-1 &\text{if $N=p-1 $},\\
0 &\text{if $p\leq N < 2p-2 $},\\
-1 &\text{if $ N = 2p-2 $},\\
z-z^p &\text{if $N=2p-1 $}.
\end{cases}
\end{equation}
\end{lemma}

\begin{lemma}\label{lem:invpol} For any prime $p$,
    if a polynomial $f(z)\in \mathbb{F}_p[z]$ satisfies $f(z+1)=f(z)$, then it is a polynomial in $z^p-z$, i.e. $f(z)\in\mathbb{F}_p[z^p-z]$.
\end{lemma}

    \begin{equation}\label{eq:polyproduct}
        \prod_{k=0}^{p-1} (z+k) = z^p - z,
    \end{equation}

\begin{lemma}\label{lem:trace}
Let $p$ be an odd prime and $M_p(z)$ be defined as in equation \eqref{eq:Mp def} with $A(z)$ a matrix polynomial of degree two,
\[ A(z) =A_0+z A_1+z^2 A_2 .\]
Then, in characteristic $p$, the trace of $M_p(z)$ takes the form,
    \begin{equation}\label{eq:trace with Ak}
    \operatorname{Tr}M_p(z) =  (\operatorname{Tr}A_2)^pZ^2+
        \operatorname{Tr}[A_1(A_1^{p-1}-A_2^{p-1})]Z+\operatorname{Tr}M_p(0),
\end{equation}
where $ Z=z^p-z$ as in equation \eqref{eq:notationps_reform}.
\end{lemma}

\begin{proof}[Proof of Lemma \ref{lem:trace}]
Define the polynomial
\[ \chi_p(z) = \operatorname{Tr}\left[A(z+p-1)\cdot A(z+p-2)\cdot\ldots\cdot A(z+1)\cdot A(z)\right] .\]
As the trace is invariant under circular shifts we can immediately deduce that $\chi_p(z+1) = \chi_p(z)$. Furthermore, as the degree of $A(z)$ is at most two, we conclude that the degree of $\chi_p(z)$ is bounded by $2p$. Thus, application of Lemma \ref{lem:invpol} yields
\begin{equation*}
    \chi_p(z)=c_0+c_1 (z^p-z)+c_2(z^p-z)^2,
\end{equation*}
for some coefficients $c_i$, $0\leq i\leq 2$, which are polynomials in the entries of the $A_k$, $0\leq k\leq 2$, over $\mathbb{F}_p$. It remains to determine $c_1$ and $c_2$.

We introduce the notation $\mathbf{j}=(j_0,\ldots,j_{p-1})\in \{0,1,2\}^p$ for length $p$ tuples of indices and define the set of indices that sum up to $d$,
  \begin{equation*}
      J_d=\{\mathbf{j}\in \{0,1,2\}^p:j_0+\ldots+j_{p-1}=d\},
  \end{equation*}
for $0\leq d\leq 2p$. We then have a cyclic action on $J_d$ generated by the permutation
\begin{equation*}
    \sigma: J_d\rightarrow J_d, \qquad \mathbf{j}\mapsto (j_1,\ldots,j_{p-1},j_0).
\end{equation*}
Fixed points of $\sigma$ are given by $\mathbf{j}'s$ with all coordinates equal,  $j_0=j_1=...=j_{p-1}$. This is only possible if $p$ divides $d$. As $p$ is a prime, every other orbit under $\sigma$ has length $p$.

We  decompose $\chi_p(z)$ as follows,
    \begin{equation}\label{eq:trace_decomp}
       \chi_p(z)=\sum_{d=0}^{2p} H_d(z),
    \end{equation}
where
\begin{equation}\label{eq:H(z) sum}
   H_d(z)=\sum_{\mathbf{j}\in J_d}g_\mathbf{j}(z) \operatorname{Tr}\left[A_{j_{p-1}}A_{j_{p-2}}\cdot\ldots\cdot A_{j_1}A_{j_0}\right],
\end{equation}
with
\begin{equation*}
    g_\mathbf{j}(z)=(z+p-1)^{j_{p-1}}(z+p-2)^{j_{p-2}}\cdot\ldots\cdot (z+1)^{j_1}(z+0)^{j_0}.
\end{equation*}
Note that by construction, the degree of $H_d(z)$ is at most $d$. This means that to compute $c_1$ and $c_2$ we only need to consider contributions from $H_d(z)$ for $p\leq d \leq 2p$ in equation \eqref{eq:trace_decomp}. We work out these contributions separately for the cases $d=2p$, $d=2p-1$, $p<d<2p-1$ and $d=p$.\\
$\mathbf{d=2p}$: We have $J_{2p}=\{(2,2,\ldots,2,2)\}$ is a singleton and
    \begin{eqnarray}
        H_{2p}(z) &=& (z+p-1)^2...(z+1)^2z^2\operatorname{Tr}A_{2}^p,\nonumber\\
        &=& [(z+p-1)...(z+1)z]^2\operatorname{Tr}A_{2}^p,\nonumber\\
        &=& (z^p-z)^2\operatorname{Tr}A_{2}^p,\label{eq:c3}
    \end{eqnarray} 
    where the last equality is a direct application of equation \eqref{eq:polyproduct}. We conclude that $c_2 = \operatorname{Tr}A_{2}^p$.\\
$\mathbf{d=2p-1}$: In this case $J_{d}$ has $p$ elements, each with all indices equal to $2$ but one (whose value is $1$), so that
    \begin{align}
        H_{2p-1}(z) &= \left( \sum_{i=0}^{p-1}\frac{\prod_{k=0}^{p-1} (z+k)^2}{z+i} \right) \operatorname{Tr}A_{2}^{p-1}A_1
        = \left(\prod_{k=0}^{p-1} (z+k)\right) \left( \sum_{i=0}^{p-1}\frac{\prod_{k=0}^{p-1} (z+k)}{z+i} \right) \operatorname{Tr}A_{2}^{p-1}A_1\nonumber\\
        &= (z^p-z)  \frac{\partial}{\partial z} \left( \prod_{k=0}^{p-1} (z+k)\right) \operatorname{Tr}A_{2}^{p-1}A_1=  (z^p-z)  \frac{\partial}{\partial z}(z^p-z) \operatorname{Tr}A_{2}^{p-1}A_1\nonumber\\
        &=  -(z^p-z)\operatorname{Tr}A_{2}^{p-1}A_1. \label{eq:c2(i)}
    \end{align} 
    The first equality follows from the fact that the trace is invariant under circular shifts. The second to last equality follows from \eqref{eq:polyproduct}.
    \item $\mathbf{p< d <2p-1}$: 
As $p$ does not divide $d$, every orbit under $\sigma$ in $J_d$ has length $p$. Let us choose a set of representatives of the corresponding quotient,
\begin{equation*}
    J_d/\langle \sigma \rangle =\{[\mathbf{j}^{(r)}]:1\leq r \leq R\},
\end{equation*}
where $R=|J_d|/p$. We have, by the cycling invariance of traces of products of matrices,
\begin{equation*}
   H_d(z)=\sum_{r=1}^R\Big( \sum_{n=0}^{p-1}g_{\mathbf{j}^{(r)}}(z+n) \Big) \operatorname{Tr}\left[A_{j_{p-1}^{(r)}}A_{j_{p-2}^{(r)}}\cdot\ldots\cdot A_{j_1^{(r)}}A_{j_0^{(r)}}\right].
\end{equation*}
By Lemma \ref{lemma:z+n}, we know that each of the inner sums
\begin{equation*}
    \sum_{n=0}^{p-1} g_{\mathbf{j}^{(r)}}(z+n),
\end{equation*}
is a constant. In particular, they do not contribute to $c_1$ or $c_2$.\\
$\mathbf{d=p}$: Following the same arguments as those used for $p<d<2p-1$, we conclude that the only term that will contribute to $c_1$ is the fixed point $\mathbf{j} = (1,1,...,1)$. Thus, using \eqref{eq:polyproduct} we conclude that
\begin{equation}
    H_{p}(z) = (z^p-z)\operatorname{Tr}A_1^p+c, \label{eq:c2(ii)}
\end{equation}
for some constant $c$.

Combining equations \eqref{eq:c3}, \eqref{eq:c2(i)} and \eqref{eq:c2(ii)}, we obtain  equation \eqref{eq:trace with Ak} and thus the lemma.
\end{proof}

\end{document}
